\chapter{Pendahuluan}

%=========================================================%
%              TULIS ISI PENDAHULUAN DI SINI              %
% Jangan Lupa untuk Menghapus Contoh Tulisan di Bawah Ini %
%          (Termasuk Panduan di Setiap Subbabnya)         %
% Contoh Teks Dapat Di-comment dengan Tanda "%" di Depan  %
%   jika Ingin Mempertahankan Panduan di Dalam File Ini   %
%=========================================================%

% --- Latar Belakang Masalah ---
\section{Latar Belakang}

Uraikan bagian ini dengan metode deduktif (dari hal umum ke khusus). Mulailah dengan menjelaskan kondisi ideal atau harapan berdasarkan teori/regulasi yang ada. Selanjutnya, paparkan fakta yang realistis di lapangan yang menunjukkan adanya kesenjangan (\textit{gap}) dengan kondisi ideal tersebut, yang mendasari mengapa penelitian ini penting untuk dilakukan. Argumen Anda dapat didukung pula dengan menyajikan data awal (grafik, tabel, atau hasil observasi ringkas) serta rujukan singkat dari penelitian terdahulu yang sejenis untuk memperkuat kepentingan posisi penelitian Anda.



% --- Rumusan Masalah ---
\section{Rumusan Masalah}

Tuliskan inti dari permasalahan yang akan Anda teliti secara terfokus, tajam, dan spesifik. Bagian ini juga dikenal sebagai pertanyaan penelitian (\textit{research questions}). Anda dapat menyusunnya dalam bentuk kalimat tanya yang membutuhkan jawaban ilmiah yang mendalam, atau berupa pernyataan tegas mengenai substansi konflik variabel yang ingin dipecahkan melalui penelitian ini.



% --- Tujuan Penelitian ---
\section{Tujuan Penelitian}

Uraikan target atau sasaran yang ingin dicapai melalui penelitian ini. Perlu diperhatikan bahwa tujuan penelitian merupakan representasi dari jawaban atas pertanyaan yang telah Anda susun di bagian Perumusan Masalah. Oleh karena itu, redaksi penulisan di bagian ini harus berkorespondensi secara langsung, runtut, dan konsisten dengan butir-butir yang ada pada perumusan masalah (misal: jika rumusan masalah ada dua, maka butir tujuan penelitian juga berjumlah dua).



% --- Manfaat Penelitian ---
\section{Manfaat Penelitian}

Uraikan secara jelas kontribusi kontemplatif dari hasil penelitian Anda jika target tujuan telah tercapai. Bagian ini umumnya dibagi menjadi dua aspek utama: (1) Manfaat Teoretis/Akademis, yaitu sumbangsih hasil penelitian bagi perkembangan khazanah ilmu pengetahuan dan masyarakat keilmuan yang relevan; dan (2) Manfaat Praktis, yaitu kegunaan konkret bagi institusi tempat penelitian, pembuat kebijakan, atau pelaku industri dalam upaya pemecahan masalah serta pengembangan kualitas kerja di lapangan.